Some retinal diseases are the result of altered genes and can lead to progressive vision loss. Usher syndrome is one such condition that impacts vision and hearing simultaneously. With recent advances in gene therapy research, it has become increasingly important to measure the impacts of therapy objectively and efficiently. This project addresses this need using a use case involving patients with Usher syndrome. Two image-based data types are used. The structural type consists of optical coherence tomography (OCT) scans, which provide cross-sectional images of the eye. The second data type is microperimetry, which measures visual sensitivity at different points of the retina. Microperimetry provides a functional measurement; however, it is time-consuming and depends on patient concentration, making it difficult to repeat reliably, especially for patients with Usher syndrome who experience visual and auditory loss simultaneously. This thesis examines whether OCT can provide useful information about retinal sensitivity, potentially supporting more practical monitoring without treating imaging as a replacement for functional testing. Previous studies have demonstrated relationships between retinal structure and function in several diseases, but there is limited evidence for this approach in small and heterogeneous Usher syndrome cohorts. The study therefore develops a transparent workflow in which matched microperimetry and OCT examinations are spatially aligned, and local structural information is compared with the sensitivity measured at each retinal point. The scope of this project is narrow, with data from 22 patients and 37 recorded MAIA sensitivity points per accepted eye-visit. Model performance is assessed using complementary measures of prediction error, agreement and calibration, with uncertainty estimated at patient level. The resulting framework provides a reproducible investigation into retinal structure--function relationships in an Usher syndrome cohort and analyses what functional predictions can reasonably be inferred from OCT structure.
